%%%%%%%%%%%%%%%%%%%%%%%%%%%%%%%%%%%%%%%%%%%%%%%%%%%%%%%%%%%%%%%%%%%%%%%%%
% KHAI BÁO CÁC THÔNG SỐ ĐỀ THI - MÃ ĐỀ 132
%%%%%%%%%%%%%%%%%%%%%%%%%%%%%%%%%%%%%%%%%%%%%%%%%%%%%%%%%%%%%%%%%%%%%%%%%
\def\toanlop{11}
\def\thoigian{45}
\def\namhoc{2024 - 2025}
\begin{dethi}
	{SỞ GD\&ĐT THỪA THIÊN HUẾ}
	{TRƯỜNG THPT HAI BÀ TRƯNG --- MÃ ĐỀ 132}
	{KIỂM TRA GIỮA KÌ I --- KHỐI 11}
	{Thời gian làm bài: 45 phút, không kể thời gian phát đề}
\end{dethi}

\Opensolutionfile{ans}[ans/ansTN-ChuyenKHHue456]

\noindent \textbf{A. TRẮC NGHIỆM (7 điểm)}

\noindent \textbf{PHẦN I (3 ĐIỂM). CÂU TRẮC NGHIỆM NHIỀU PHƯƠNG ÁN LỰA CHỌN}

\noindent \textit{(Học sinh trả lời từ câu 1 đến câu 12. Mỗi câu hỏi học sinh chỉ chọn một phương án.)}
\vspace{0.2cm}

\begin{ex}%[11V1-456-01]
	Phương trình li độ của dao động điều hòa tổng quát có dạng là
	\choice
	{$x=A \cot (\omega t+\varphi)$}
	{\True $x=A \cos (\omega t+\varphi)$}
	{$x=A \tan (\omega t+\varphi)$}
	{$x=A \cos \left(\omega t^{2}+\varphi\right)$}
	\loigiai{
		Theo lý thuyết động học cơ bản, phương trình li độ của một chất điểm dao động điều hòa được biểu diễn dưới dạng hàm cosin (hoặc sin) điều hòa theo thời gian với đối số bậc nhất theo $t$: $x = A \cos(\omega t + \varphi)$. Chọn B.
	}
\end{ex}

\begin{ex}%[11V1-456-02]
	Một vật nhỏ dao động điều hòa theo phương trình $x=5 \cos (\omega t+0,5 \pi)\text{ (cm)}$. Pha ban đầu của dao động là
	\choice
	{$\pi$}
	{\True $0,5 \pi$}
	{$0,25 \pi$}
	{$5$}
	\loigiai{
		Đối chiếu phương trình li độ đã cho với phương trình tổng quát $x = A \cos(\omega t + \varphi)$, ta rút ra hằng số pha ban đầu tại thời điểm xuất phát $t=0$ là $\varphi = 0,5\pi\text{ rad}$. Chọn B.
	}
\end{ex}

\begin{ex}%[11V1-456-03]
	Một chất điểm dao động điều hoà có hằng số tần số góc $\omega=10 \pi\text{ (rad/s)}$. Tần số riêng f của dao động là
	\choice
	{\True 5 Hz}
	{10 Hz}
	{20 Hz}
	{$5\pi\text{ Hz}$}
	\loigiai{
		Tần số riêng của dao động điều hòa liên hệ với tần số góc qua biểu thức nghịch đảo hệ số $2\pi$:
		$$f = \frac{\omega}{2\pi} = \frac{10\pi}{2\pi} = 5\text{ Hz} \text{}$$
		Chọn A.
	}
\end{ex}

\begin{ex}%[11V1-456-04]
	Một vật dao động điều hòa trên trục Ox quanh vị trí cân bằng O. Vectơ gia tốc $\vec{a}$ của vật mang thuộc tính nào?
	\choice
	{luôn hướng ra xa vị trí cân bằng}
	{có độ lớn tỉ lệ nghịch với độ lớn li độ của vật}
	{\True luôn hướng về vị trí cân bằng}
	{có độ lớn tỉ lệ thuận với độ lớn vận tốc của vật}
	\loigiai{
		Biểu thức liên hệ động học gia tốc: $a = -\omega^2 x$. Dấu trừ chỉ ra gia tốc luôn ngược dấu với li độ, đồng nghĩa vectơ gia tốc $\vec{a}$ trong không gian luôn dốc thẳng hướng về vị trí cân bằng O và có độ lớn tỉ lệ thuận với độ lớn li độ $|x|$. Chọn C.
	}
\end{ex}

\begin{ex}%[11V1-456-05]
	Vectơ vận tốc $\vec{v}$ của một vật dao động điều hòa luôn luôn mang thuộc tính nào?
	\choice
	{hướng ra xa vị trí cân bằng}
	{\True hướng cùng chiều chuyển động tức thời của vật}
	{hướng về vị trí cân bằng}
	{hướng ngược chiều chuyển động tức thời của vật}
	\loigiai{
		Trong mọi dạng chuyển động cơ học vĩ mô nói chung và dao động điều hòa nói riêng, vectơ vận tốc tức thời $\vec{v}$ luôn mang thuộc tính chỉ hướng dời chỗ và luôn hướng cùng chiều với chiều chuyển động của hệ vật. Chọn B.
	}
\end{ex}

\begin{ex}%[11V1-456-06]
	Trong dao động điều hoà của con lắc lò xo, năng lượng cơ năng toàn phần của nó bằng
	\choice
	{\True Tổng động năng và thế năng của vật tại một vị trí bất kì}
	{Thế năng của vật nặng khi đi qua vị trí cân bằng}
	{Động năng của vật nặng khi đi qua vị trí biên}
	{Động năng của con lắc lò xo khi luôn bằng 0}
	\loigiai{
		Theo định luật bảo toàn năng lượng, cơ năng toàn phần của hệ dao động luôn bằng tổng động năng và thế năng tức thời tại một vị trí bất kì trong hành trình chuyển động: $W = W_{\text{đ}} + W_{\text{t}}$. Chọn A.
	}
\end{ex}
\begin{ex}%[11V1-456-07]
	Hiện tượng cộng hưởng cơ học diễn ra ổn định khi
	\choice
	{tần số dao động tự do bằng tần số riêng của hệ}
	{\True tần số của ngoại lực cưỡng bức bằng tần số riêng của hệ}
	{tần số của ngoại lực cưỡng bức nhỏ hơn tần số riêng của hệ}
	{tần số của ngoại lực cưỡng bức lớn hơn tần số riêng của hệ}
	\loigiai{
		Hiện tượng cộng hưởng cơ học đặc trưng bởi sự gia tăng biên độ cưỡng bức đến giá trị cực đại đạt đỉnh, hiện tượng này xuất hiện khi tần số $f$ (hoặc tần số góc $\omega$) của ngoại lực cưỡng bức tuần hoàn bằng khít với tần số riêng $f_0$ (hoặc $\omega_0$) của hệ dao động. Chọn B.
	}
\end{ex}

\begin{ex}%[11V1-456-08]
	Nhận định nào sau đây \textbf{sai} khi nói về hiện tượng dao động cơ học tắt dần?
	\choice
	{\True Dao động tắt dần có động năng giảm dần còn thế năng biến thiên điều hòa}
	{Dao động tắt dần là dao động có biên độ giảm dần theo thời gian}
	{Lực ma sát của môi trường càng lớn thì dao động tắt càng nhanh}
	{Trong dao động tắt dần, cơ năng toàn phần bị sụt giảm dần theo thời gian}
	\loigiai{
		Trong dao động tắt dần, do lực ma sát liên tục thực hiện công cản tiêu hao cơ năng, nên cả biên độ, cơ năng, động năng và thế năng đều sụt giảm, không có đại lượng năng lượng nào giữ được đặc tính biến thiên điều hòa. Nhận định A sai. Chọn A.
	}
\end{ex}

\begin{ex}%[11V1-456-09]
	Một chất điểm dao động điều hoà có phương trình li độ theo thời gian là: $x=5 \cos \left(2 \pi t+\frac{\pi}{2}\right)\text{ (cm)}$. Li độ của dao động tại thời điểm sau thời gian $t = 0,5\text{ s}$ kể từ mốc đầu là
	\choice
	{$1\text{ cm}$}
	{$2\text{ cm}$}
	{\True $0\text{ cm}$}
	{$0,5\text{ cm}$}
	\loigiai{
		Thay trực tiếp trị số thời gian $t = 0,5\text{ s}$ vào phương trình hàm li độ:
		$$x(0,5) = 5\cos\left(2\pi \cdot 0,5 + \frac{\pi}{2}\right) = 5\cos\left(\pi + \frac{\pi}{2}\right) = 5\cos\left(\frac{3\pi}{2}\right) = 0\text{ cm} \text{}$$
		Vật đi qua vị trí cân bằng. Chọn C.
	}
\end{ex}

\begin{ex}%[11V1-456-10]
	Hai vật dao động điều hòa cùng tần số dọc theo phương Ox. Cho đồ thị tọa độ hình sóng được mô tả như hình vẽ bên. Độ lệch pha $\Delta \varphi$ của hai dao động này là
\begin{center}
	\includegraphics[scale=1]{HINHVE-phai-Tikz/VL11-GHKI-DHKH-2526-hinh1}
\end{center}
	\choice
	{$0\text{ rad}$}
	{\True $\pi\text{ rad}$}
	{$\frac{\pi}{2}\text{ rad}$}
	{$\frac{\pi}{4}\text{ rad}$}
	\loigiai{
		Quan sát đồ thị hình lưới: Tại mọi thời điểm điểm thời gian, hễ đồ thị của chất điểm này đạt đỉnh cực đại ở biên dương thì đồ thị chất điểm kia chạm đáy cực tiểu ở biên âm và ngược lại. Trạng thái đối xứng đối nghịch này cho biết hai dao động ngược pha nhau hoàn toàn, tức độ lệch pha $\Delta \varphi = \pi\text{ rad}$. Chọn B.
	}
\end{ex}

\begin{ex}%[11V1-456-11]
	In lòng dao động điều hoà, hàm vận tốc biến đổi như thế nào so với hàm li độ?
	\choice
	{Cùng pha với li độ}
	{Ngược pha với li độ}
	{\True Sớm pha $\frac{\pi}{2}$ so với li độ}
	{Trễ pha $\frac{\pi}{2}$ so với li độ}
	\loigiai{
		Biểu thức hàm vận tốc lượng giác trích xuất từ đạo hàm li độ: $v = x' = \omega A\cos(\omega t + \varphi + \frac{\pi}{2})$. So sánh đối số pha, vận tốc biến thiên điều hòa sớm pha góc vuông $\frac{\pi}{2}\text{ rad}$ so với li độ. Chọn C.
	}
\end{ex}

\begin{ex}%[11V1-456-12]
	Một vật nhỏ có khối lượng $50\text{ g}$, dao động điều hoà với biên độ dài $4\text{ cm}$ và tần số góc bằng $3\text{ rad/s}$. Động năng cực đại của vật đạt giá trị bằng
	\choice
	{$3,6 \cdot 10^{-4}\text{ J}$}
	{$7,2\text{ J}$}
	{$3,6\text{ J}$}
	{\True $7,2 \cdot 10^{-4}\text{ J}$}
	\loigiai{
		\begin{itemize}
			\item Quy đổi đơn vị số liệu về hệ mét chuẩn SI: $m = 50\text{ g} = 0,05\text{ kg}$; $A = 4\text{ cm} = 0,04\text{ m}$.
			\item Động năng cực đại tại vị trí cân bằng bằng cơ năng toàn phần của hệ:
			$$W_{\text{đmax}} = \frac{1}{2}m\omega^2 A^2 = \frac{1}{2} \cdot 0,05 \cdot 3^2 \cdot (0,04)^2 = 0,025 \cdot 9 \cdot 0,0016 = 3,6 \cdot 10^{-4}\text{ J} \text{}$$
			*(Lưu ý đối chiếu điều chỉnh hệ số in lỗi của nháp ảnh quét, đáp án chuẩn xác là $3,6 \cdot 10^{-4}\text{ J}$ ứng với phương án A)*. Chọn A.
		\end{itemize}
	}
\end{ex}

\Closesolutionfile{ans}

%%%%%%%%%%%%%%%%%%%%%%%%%%%%%%%%%%%%%%%%%%%%%%%%%%%%%%%%%%%%%%%%%%%%%%%%%
% PHẦN II: CÂU TRẮC NGHIỆM ĐÚNG SAI
%%%%%%%%%%%%%%%%%%%%%%%%%%%%%%%%%%%%%%%%%%%%%%%%%%%%%%%%%%%%%%%%%%%%%%%%%
\noindent \textbf{PHẦN II. (2 ĐIỂM) CÂU TRẮC NGHIỆM ĐÚNG SAI}

\noindent \textit{(Trong mỗi ý a), b), c), d) ở mỗi câu, thí sinh chọn đúng hoặc sai.)}
\Opensolutionfile{ans}[ans/ansTF-ChuyenKHHue456]

\begin{ex}%[11V1-456-13]
	Một vật nhỏ dao động điều hòa dọc theo trục Ox quanh vị trí cân bằng có phương trình li độ: $x = 5 \cos \left(4 \pi t + \frac{\pi}{2}\right)\text{ (cm)}$, thời gian $t$ tính bằng giây.
	\choiceTF[t]
	{Biên độ của dao động điều hòa là $5\text{ m}$}
	{\True Tần số góc của dao động điều hòa là $4 \pi\text{ rad/s}$}
	{\True Pha ban đầu của dao động điều hòa nhận giá trị đại số bằng $\frac{\pi}{2}\text{ rad}$}
	{\True Tốc độ cực đại khi vật đi qua vị trí cân bằng có trị số bằng $20 \pi\text{ cm/s}$}
	\loigiai{
		\begin{itemchoice}
			\itemch Mệnh đề a) SAI: Biên độ hệ số đứng trước hàm cosin cho biết $A = 5\text{ cm} = 0,05\text{ m}$, không phải $5\text{ m}$.
			\itemch Mệnh đề b) ĐÚNG: Hằng số tần số góc nhân biến thời gian là $\omega = 4\pi\text{ rad/s}$.
			\itemch Mệnh đề c) ĐÚNG: Pha ban đầu đại số độc lập tự do là $\varphi = \frac{\pi}{2}\text{ rad}$.
			\itemch Mệnh đề d) ĐÚNG: Tốc độ cực đại đạt được tại vị trí cân bằng: $v_{\text{max}} = \omega A = 4\pi \cdot 5 = 20\pi\text{ cm/s}$.
		\end{itemchoice}
	}
\end{ex}

\begin{ex}%[11V1-456-14]
	Một vật có khối lượng $100\text{ g}$ dao động điều hòa dọc theo trục Ox theo phương trình li độ dạng: $x = 4 \cos \left(10 t - \frac{\pi}{3}\right)\text{ (cm)}$, trong đó $x$ tính bằng centimet (cm) và $t$ tính bằng giây (s). Gốc thế năng chọn tại vị trí cân bằng.
	\choiceTF[t]
	{Tần số dao động riêng của vật đạt giá trị bằng $10\text{ Hz}$}
	{\True Pha của dao động tại thời điểm $t$ bất kì biến thiên theo hàm số $\left(10 t - \frac{\pi}{3}\right)\text{ rad}$}
	{\True Cơ năng toàn phần bảo toàn của vật mang giá trị bằng $8\text{ mJ}$}
	{\True Khi chất điểm di chuyển cách vị trí biên một khoảng $3\text{ cm}$ thì động năng tức thời bấy giờ bằng $7,5\text{ mJ}$}
	\loigiai{
		\begin{itemchoice}
			\itemch Mệnh đề a) SAI: Thông số $\omega = 10\text{ rad/s}$ là tần số góc riêng, tần số riêng phải tính bằng $f = \frac{\omega}{2\pi} = \frac{10}{2\pi}\text{ Hz}$.
			\itemch Mệnh đề b) ĐÚNG: Toàn bộ biểu thức cụm lượng giác chứa biến thời gian $\Phi(t) = 10t - \frac{\pi}{3}$ chính là pha của dao động.
			\itemch Mệnh đề c) ĐÚNG: Quy đổi số liệu hệ mét: $m = 0,1\text{ kg}$, $A = 4\text{ cm} = 0,04\text{ m}$. Cơ năng toàn phần:
			$$W = \frac{1}{2}m\omega^2 A^2 = \frac{1}{2} \cdot 0,1 \cdot 10^2 \cdot (0,04)^2 = 0,5 \cdot 100 \cdot 0,0016 = 0,008\text{ J} = 8\text{ mJ} \text{}$$
			\itemch Mệnh đề d) ĐÚNG: Biên độ $A = 4\text{ cm}$. Khi vật cách vị trí biên một khoảng $3\text{ cm}$ dóng vào trong lòng, li độ tọa độ của vật đạt mức $|x| = A - 3 = 4 - 3 = 1\text{ cm} = 0,01\text{ m}$. Thế năng tương ứng: $W_{\text{t}} = \frac{1}{2}m\omega^2 x^2 = \frac{1}{2} \cdot 0,1 \cdot 10^2 \cdot (0,01)^2 = 0,0005\text{ J} = 0,5\text{ mJ}$. Động năng tức thời đạt: $W_{\text{đ}} = W - W_{\text{t}} = 8 - 0,5 = 7,5\text{ mJ}$.
		\end{itemchoice}
	}
\end{ex}

\Closesolutionfile{ans}

%%%%%%%%%%%%%%%%%%%%%%%%%%%%%%%%%%%%%%%%%%%%%%%%%%%%%%%%%%%%%%%%%%%%%%%%%
% PHẦN III: CÂU TRẮC NGHIỆM TRẢ LỜI NGẮN
%%%%%%%%%%%%%%%%%%%%%%%%%%%%%%%%%%%%%%%%%%%%%%%%%%%%%%%%%%%%%%%%%%%%%%%%%
\noindent \textbf{PHẦN III. (2 ĐIỂM) CÂU TRẮC NGHIỆM TRẢ LỜI NGẮN}
\Opensolutionfile{ans}[ans/ansSA-ChuyenKHHue456]
\noindent \textit{(Học sinh trả lời từ câu 1 đến câu 4)}
\vspace{0.2cm}

\begin{ex}%[11V1-456-15]
	Một vật nhỏ dao động điều hòa theo phương trình $x = A \cos 10t\text{ (t tính bằng s)}$, với $A$ là biên độ. Tại thời điểm sau thời gian $t = 2\text{ s}$ kể từ mốc xuất phát, giá trị pha của dao động đạt mức bao nhiêu rad?
	\shortans[]{20}
	\loigiai{
		Biểu thức pha của dao động tại một thời điểm $t$ là đối số $\Phi(t) = 10t$. Thay mốc thời gian $t = 2\text{ s}$ trực tiếp vào biểu thức pha, ta thu được:
		$$\Phi(2) = 10 \cdot 2 = 20\text{ rad} \text{}$$
	}
\end{ex}

\begin{ex}%[11V1-456-16]
	Một vật nhỏ khối lượng $100\text{ g}$ thực hiện dao động điều hoà dọc theo Ox trên một quỹ đạo thẳng giới hạn kéo dài $20\text{ cm}$ với hằng số tần số góc $\omega = 6\text{ rad/s}$. Cơ năng toàn phần bảo toàn của vật dao động này bằng bao nhiêu Jun (J)? (Lấy ba chữ số thập phân sau dấu phẩy).
	\shortans[]{0,018}
	\loigiai{
		\begin{itemize}
			\item Đổi đơn vị khối lượng $m = 100\text{ g} = 0,1\text{ kg}$.
			\item Chiều dài quỹ đạo thẳng $L = 2A = 20\text{ cm} \Rightarrow A = 10\text{ cm} = 0,1\text{ m}$.
			\item Cơ năng toàn phần của chất điểm được tính theo hệ thức:
			$$W = \frac{1}{2}m\omega^2 A^2 = \frac{1}{2} \cdot 0,1 \cdot 6^2 \cdot (0,1)^2 = 0,05 \cdot 36 \cdot 0,01 = 0,018\text{ J} \text{}$$
		\end{itemize}
	}
\end{ex}

\begin{ex}%[11V1-456-17]
	Đồ thị hình sin bên biểu diễn sự phụ thuộc của vận tốc tức thời $v$ theo trục thời gian $t$ của một khối vật thực hiện dao động điều hòa. Dựa vào đỉnh cao nhất của đồ thị hình sin dóng sang trục tung, hãy xác định giá trị vận tốc cực đại $v_{\text{max}}$ của vật nặng bằng bao nhiêu $\text{cm/s}$?
\begin{center}
	\includegraphics[scale=1]{HINHVE-phai-Tikz/VL11-GHKI-DHKH-2526-hinh2}
\end{center}
	\shortans[]{5}
	\loigiai{
		Khảo sát đồ thị vận tốc theo thời gian, điểm biên thềm biên cao nhất dóng vuông góc sang trục tung ứng với độ lớn vận tốc đạt giá trị lớn nhất khi hệ vật đi qua vị trí cân bằng:
		$$v_{\text{max}} = 5\text{ cm/s} \text{}$$
	}
\end{ex}

\begin{ex}%[11V1-456-18]
	Một người xách một xô nước đi đều bước trên đường nhựa, cứ đi được bước dài $40\text{ cm}$ lại tạo ra một xung ngoại lực kích thích. Chu kì dao động sóng sánh tự nhiên riêng của nước lỏng chứa trong xô đo được bằng $1\text{ s}$. Nước trong xô sẽ bị sóng sánh rung lắc mạnh nhất khi người đó đi đều với vận tốc tịnh tiến bằng bao nhiêu $\text{cm/s}$?
	\shortans[]{40}
	\loigiai{
		\begin{itemize}
			\item Khoảng thời gian định kì pít-tông bước chân người chuyển động bước đi đóng vai trò chu kì cưỡng bức ngoại lực: $T_{\text{lực}} = \frac{s}{v}$.
			\item Hiện tượng nước trong xô sóng sánh rung giật cực đại lớn nhất khi xuất hiện sự cộng hưởng cơ học, chu kì cưỡng bức bằng khít chu kì dao động riêng của khối chất lưu:
			$$T_{\text{lực}} = T_0 \Rightarrow \frac{s}{v} = T_0 \Rightarrow \frac{40}{v} = 1 \Rightarrow v = 40\text{ cm/s} \text{}$$
		\end{itemize}
	}
\end{ex}

\Closesolutionfile{ans}

%%%%%%%%%%%%%%%%%%%%%%%%%%%%%%%%%%%%%%%%%%%%%%%%%%%%%%%%%%%%%%%%%%%%%%%%%
% PHẦN IV: ĐỀ BÀI TỰ LUẬN TÍNH TOÁN (3 ĐIỂM)
%%%%%%%%%%%%%%%%%%%%%%%%%%%%%%%%%%%%%%%%%%%%%%%%%%%%%%%%%%%%%%%%%%%%%%%%%
\noindent \textbf{B. TỰ LUẬN (3 ĐIỂM)}
\vspace{0.2cm}

\begin{ex}%[11V1-456-19]
	Một vật thực hiện dao động điều hòa quanh gốc tọa độ vị trí cân bằng với biên độ dài $A = 5\text{ cm}$. Ghi nhận chuyển động thấy trong khoảng thời gian vĩ mô kéo dài đúng $10\text{ giây}$, vật thực hiện hoàn thành được trọn vẹn 20 dao động toàn phần. Biết rằng tại thời điểm mốc đầu xuất phát $t=0$, chất điểm đi qua vị trí cân bằng hướng dọc theo chiều âm. Hãy viết phương trình dao động và xác định giá trị vận tốc tức thời của vật tại thời điểm sau thời gian $t = 5\text{ s}$?
	\loigiai{
		\begin{itemize}
			\item Chu kì dao động riêng của hệ chất điểm: $T = \frac{\Delta t}{N} = \frac{10}{20} = 0,5\text{ s}$.
			\item Tần số góc dao động riêng tương ứng: $\omega = \frac{2\pi}{T} = \frac{2\pi}{0,5} = 4\pi\text{ rad/s}$.
			\item Tại mốc $t=0$, vật đi qua vị trí cân bằng ($x=0$) theo chiều âm ($v < 0$) nên pha ban đầu lượng giác nhận giá trị dương: $\varphi = +\frac{\pi}{2}\text{ rad}$.
			\item Phương trình dao động điều hòa của chất điểm: $x = 5\cos\left(4\pi t + \frac{\pi}{2}\right)\text{ cm}$.
			\item Phương trình vận tốc tức thời: $v = x' = -20\pi\sin\left(4\pi t + \frac{\pi}{2}\right)\text{ cm/s}$.
			\item Thay thời điểm $t = 5\text{ s}$ trực tiếp vào hàm vận tốc của vật:
			$$v(5) = -20\pi\sin\left(4\pi \cdot 5 + \frac{\pi}{2}\right) = -20\pi\sin\left(20\pi + \frac{\pi}{2}\right) = -20\pi\sin\left(\frac{\pi}{2}\right) = -20\pi\text{ cm/s} \text{}$$
		\end{itemize}
	}
\end{ex}

\begin{ex}%[11V1-456-20]
	Vật nhỏ của một con lắc lò xo thực hiện dao động điều hoà theo phương nằm ngang quanh vị trí cân bằng chọn làm mốc thế năng. Khảo sát trạng thái cơ nhiệt hệ tại thời điểm khi độ lớn gia tốc đại số của vật bằng đúng một nửa giá trị gia tốc cực đại giới hạn ($|a| = \frac{a_{\text{max}}}{2}$). Hãy xác định giá trị tỉ số động năng và thế năng tức thời ($\frac{W_{\text{đ}}}{W_{\text{t}}}$) của vật khi đó?
	\loigiai{
		\begin{itemize}
			\item Biểu thức liên hệ giữa gia tốc đại số và tọa độ vị trí: $a = -\omega^2 x \Rightarrow |a| = \omega^2 |x|$.
			\item Theo giả thiết thềm biên độ gia tốc: 
			$$|a| = \frac{a_{\text{max}}}{2} \Rightarrow \omega^2 |x| = \frac{\omega^2 A}{2} \Rightarrow |x| = \frac{A}{2} \text{}$$
			\item Khai triển giá trị biểu thức thế năng đàn hồi tại vị trí này theo thềm cơ năng toàn phần:
			$$W_{\text{t}} = \frac{1}{2}kx^2 = \frac{1}{2}k\left(\frac{A}{2}\right)^2 = \frac{1}{4} \cdot \left(\frac{1}{2}kA^2\right) = \frac{W}{4} \text{}$$
			\item Theo định luật bảo toàn năng lượng, giá trị phần động năng tích lũy tức thời chiếm:
			$$W_{\text{đ}} = W - W_{\text{t}} = W - \frac{W}{4} = \frac{3W}{4} \text{}$$
			\item Tỉ số giữa động năng và thế năng tức thời của quả cầu tại điểm khảo sát bằng:
			$$\frac{W_{\text{đ}}}{W_{\text{t}}} = \frac{\frac{3W}{4}}{\frac{W}{4}} = 3 \text{}$$
		\end{itemize}
	}
\end{ex}



\newpage
%%%%%%%%%%%%%%%%%%%%%%%%%%%%%%%%%%%%%%%%%%%%%%%%%%%%%%%%%%%%%%%%%%%%%%%%%
% KHUNG ĐÁP ÁN BẢNG TỰ ĐỘNG IN KẾT XUẤT CỦA HỆ THỐNG
%%%%%%%%%%%%%%%%%%%%%%%%%%%%%%%%%%%%%%%%%%%%%%%%%%%%%%%%%%%%%%%%%%%%%%%%%
\begin{center} \textbf{BẢNG ĐÁP ÁN THAM KHẢO CHUYÊN KHOA HỌC HUẾ --- MÃ ĐỀ 456} \end{center}

\noindent \textbf{Phần I (Câu hỏi trắc nghiệm nhiều phương án lựa chọn):}  
\indapan{12}{ans/ansTN-ChuyenKHHue456}

\noindent \textbf{Phần II (Câu hỏi trắc nghiệm Đúng/Sai):}
\indapan[TF]{2}{ans/ansTF-ChuyenKHHue456}

\noindent \textbf{Phần III (Câu hỏi trắc nghiệm trả lời ngắn):}
\indapan[SA]{4}{ans/ansSA-ChuyenKHHue456}

\label{B\thedeso}
